\section{Broader Impact and Ethics}
\label{sec:app-ethics}

\paragraph{Dual-use considerations.}
Neural cryptanalysis research inherently carries dual-use risk: techniques that evaluate cipher security can also be used to attack deployed systems.
We mitigate this in three ways.
First, our results apply only to round-reduced variants---all three target ciphers use substantially more rounds than our distinguishers penetrate (e.g., 22 rounds for \speck vs.\ our 7-round limit).
Second, the transfer polarity finding is primarily diagnostic: it reveals \emph{what} neural networks learn, not how to break full ciphers.
Third, we report all negative results (failed key recovery on \simon, failed RL search) alongside positive ones, providing an honest assessment of the method's limitations.

\paragraph{Responsible disclosure.}
Our work does not reveal any new vulnerability in full-round implementations of \speck, \simon, or \present.
The round-reduced attacks are consistent with the known security margins established by the cipher designers and subsequent cryptanalytic literature.

\paragraph{Deployment context.}
Lightweight ciphers are deployed in resource-constrained environments (IoT, RFID, medical devices) where security margins are necessarily tighter than in general-purpose ciphers like AES.
Our cross-family comparison provides empirical evidence to inform cipher selection decisions: the vulnerability ordering (\simon $\gg$ \speck $>$ \present) may be relevant when choosing between these ciphers for constrained applications, alongside other factors such as hardware cost, latency, and energy consumption.
